\documentclass[11pt,letterpaper]{article}
\usepackage[letterpaper,top=0.75in,bottom=0.75in,left=1in,right=1in]{geometry}
\usepackage[T1]{fontenc}
\usepackage{mathptmx}
\usepackage{microtype}
\usepackage[hidelinks]{hyperref}
\usepackage{titlesec}
\setlength{\parindent}{0pt}
\setlength{\parskip}{3pt}
\titleformat{\section}{\normalfont\normalsize\bfseries}{}{0pt}{}
\titlespacing*{\section}{0pt}{6pt}{2pt}
\newcommand{\rg}[1]{\par\smallskip{\normalfont\bfseries\itshape #1}\par}
\newcommand{\task}[1]{\textbf{#1}\hspace{0.4em}}
\hypersetup{pdftitle={Trustworthy Collective Reasoning and Evaluation for AI Agents},pdfauthor={Xinya Du}}
\begin{document}
\begin{center}
{\fontsize{14}{16}\selectfont\bfseries Trustworthy Collective Reasoning and Evaluation for AI Agents\par}
\vspace{2pt}
{\small Xinya Du | Assistant Professor of Computer Science | The University of Texas at Dallas | xinya.du@utdallas.edu\par}
{\small ONR FY27 YIP | Program Officer: Dr.\ Predrag Neskovic, Code 311, Mathematical Data Science\par}
\end{center}

\section*{1\quad The problem}
Five AI agents read a set of shipping reports and are asked whether a cargo vessel has left port. Three answer yes and two answer no. The vessel is in fact still in port. Two of the yes answers rest on the same outdated schedule, and the third comes from an agent that has been compromised and argues confidently for yes. Nobody knows in advance which agents are reliable on questions like this. A majority vote picks the wrong answer, and an open discussion may do no better, because the compromised agent is persuasive and the two agents with the same source reinforce each other.

\textbf{The central problem of this project is how a group of AI agents of unknown expertise, some possibly dishonest, should weight and combine their judgments, resolve their disagreements, and improve from the outcome.}

Collective judgment fails in three ways, all present in the example. Counting every vote equally lets less reliable agents outvote experts, and agents that make the same mistakes count more than once. Agents may behave strategically: they can rate their own answers too highly, rate rivals low to raise their own standing, argue for an answer they do not hold, or coordinate with one another. In discussion, agents give way to whoever is most confident or persistent, with or without evidence. Research goal 1 (RG1) learns how much weight each agent deserves, RG2 decides which claims to check and settles disputes with evidence, and RG3 uses verified outcomes to improve both the weights and the agents. We focus on factual claims and reasoning steps that can be checked, at least in part.

\section*{2\quad Why it matters}
Using one model to judge another is now standard practice in AI evaluation and oversight \cite{zheng2023judging,verga2024juries}, and multi-agent systems increasingly decide by vote or discussion. Agreement among such models is weaker evidence than it appears. Their errors are strongly correlated \cite{kim2025correlated}, and more capable models make more similar mistakes \cite{goel2025great}. When $n$ judges have average error correlation $\rho$, they carry roughly the information of $n/(1+(n-1)\rho)$ independent judges \cite{hossain2026agreement}. In one measurement, nine judges were worth about 2.2 independent judges, items on which all nine agreed were still wrong 9.1\% of the time, and established aggregation methods closed at most 11\% of the gap to an independent panel, even with access to the correct answers \cite{kohli2026nine}. Discussion erodes independence further: one round of debate raised a panel's unanimity from 39.5\% to 95.2\% while its accuracy moved by less than a point \cite{shu2026unanimity}. When two agents checked each other's work over many rounds, collusion emerged in 94\% of runs across ten models, and stronger models within a family colluded sooner \cite{shi2026collusion}.

Naval decisions raise each of these risks. AI assistants and automated analyses differ in reliability across tasks, and that reliability is rarely known in advance. Ground truth arrives late and sparsely, if at all. An adversary may compromise an agent or inject false reports. Checking a claim can be costly, or impossible under limited connectivity. A commander or analyst still has to decide whether an AI consensus deserves trust.

\section*{3\quad How it is solved today, and the limitations}
Statistical methods for crowdsourcing estimate each rater's accuracy without an answer key, jointly with item difficulty \cite{dawid1979maximum,whitehill2009whose}, and truth-discovery methods detect sources that copy one another \cite{dong2009dependence}. They assume that raters err independently or that their dependence has a simple structure \cite{jaffe2016dependent}, and they treat each judgment as fixed. Two recent preprints use small trusted sets with language-model judges: one estimates accuracy and shared mistakes to choose a voting rule \cite{hossain2026agreement}, and the other drops redundant judges and routes questions to specialists \cite{zhu2026stopping}. Both treat judges as honest signals that never interact.

Computational social choice offers mechanisms in which a participant's ratings of others cannot change its own chance of selection \cite{kurokawa2015impartial,alon2011sum}, along with fair, manipulation-resistant panel selection \cite{flanigan2021fair,flanigan2024manipulation} and robust aggregation of AI evaluations \cite{khalaf2026maximal}. Lev et al.\ showed that an impartial mechanism can weight reviewers by reliability \cite{lev2023peernomination}. Their weights come from agreement with co-reviewers, which rewards siding with the majority; by the authors' analysis, accurate reviewers can no longer be identified once they are outnumbered. These mechanisms take judgments as given and do not check the claims being judged.

Multi-agent debate lets models critique one another \cite{du2024debate}, but majority voting accounts for most of its gains \cite{choi2025debate}, models abandon correct answers under social pressure \cite{sharma2024sycophancy,laban2023flipflop}, and a single persuasive adversary can cut a debating group's accuracy by roughly 10\% to 40\% \cite{amayuelas2024attack}. In AI safety debate, two models argue before a judge \cite{irving2018debate}, and restricting them to verified quotations helps the judge reach the truth \cite{khan2024debating}. In a recent collaborative variant, two disagreeing models locate the crux of their disagreement before a weak judge decides. Judge accuracy rose to 62.1\% from 49.2\% under debate, but the judges accepted the models' consensus almost every time, including when both models were wrong, and the method assumes two honest models \cite{jiang2026disagreement}. Colluding agents can also mislead judges when every fragment they cite is true \cite{hu2026lying}.

Peer prediction rewards an agent for reports that help predict other agents' reports, which can make honesty each agent's best strategy without ground truth. A recent extension to language models resists a single deceiver and leaves collusion open \cite{qiu2026truthfulness}. AI control tests oversight protocols against a red team that tries to subvert them, using a trusted weaker model and a small audit budget \cite{greenblatt2024control}. None of these lines of work learns whom to trust from verified outcomes while protecting the group's decision from persuasion and coordinated dishonesty.

\section*{4\quad Our approach: what is new, and why it will succeed}
We anchor every step in a small stream of verified checks that agents cannot predict. Each agent first answers on its own, splits its reasoning into checkable claims, and cites its sources; this sealed first answer is kept for later comparison. The goal is accurate, calibrated decisions within a checking budget that hold up when some agents are dishonest, and that report a claim as unresolved when the evidence runs out.

\textbf{What is new.} (i) Weights come from audited correctness and measured error dependence, so an agent can raise its weight only by being right. (ii) Skill-weighted aggregation stays impartial, and we measure how much a coalition can still gain. (iii) Checks are allocated so that a coalition that knows the policy cannot tell which claims are safe to corrupt. (iv) Disputes among many agents of correlated reliability are settled only by verifiable evidence, and each agent is rewarded for the verified evidence it contributes to the group's answer.

\textbf{Why it will succeed.} The PI's group built the main components. PRD introduced peer ranking and peer discussion among language-model evaluators and agreed better with human judgments than using the strongest single model as judge, which favors its own answers \cite{li2024prd}; our PeerEval leaderboard (peereval-ai.github.io) applies peer evaluation to many current models. FG-PRM detects six types of reasoning errors at the level of individual steps \cite{li2025fgprm}. HiPRAG trains search agents with step-level rewards to search only when needed, reaching 65.4\% and 67.2\% average accuracy on seven question-answering benchmarks with 3B and 7B models while cutting unnecessary searches to 2.3\% \cite{wu2026hiprag}. In the Amazon Nova AI Challenge, the PI co-led one of five red teams, chosen with the other competing teams from more than 90 proposals, that probed code-generating models for flaws. The PI also holds a 2024 NSF CAREER award and a 2023 Amazon Research Award on process-guided fine-tuning, and as a postdoctoral researcher led system integration for the RESIN-11 system in DARPA KAIROS.

\rg{RG1. Learn whom to trust}
\task{Task 1.1: Expertise- and dependence-aware weighting.} We will estimate each agent's skill on each type of question and each question's difficulty, extending the model of Dawid and Skene \cite{dawid1979maximum,whitehill2009whose} with shared error factors. Two agents share an error factor when they fail together on audited claims more often than their separate accuracies predict. Shared sources, the same underlying model, and overlapping reasoning steps help predict such factors \cite{dong2009dependence,goel2025great}, but agents that share a source can still reason differently, so the estimate itself comes from coinciding failures. Recent audit-set methods choose among fixed voting rules or prune redundant judges \cite{hossain2026agreement,zhu2026stopping}. Our model produces weights for each question type. Because a strategic agent may behave well only when it expects a check, the model learns only from audits the agent cannot predict (Task 2.1). We will measure how many audited claims are needed to recover the skill and dependence planted in our testbeds. In the vessel example, the two agents that relied on the same schedule fail together on audited claims, so their agreement counts for little more than one vote.

\task{Task 1.2: Impartial aggregation that resists coalitions.} When agents rate one another, as on our PeerEval leaderboard, weighting by skill can conflict with impartiality. In PRD's ranking, an agent that lowers a rival's score also lowers that rival's weight as a judge, which can raise its own score. We will take each agent's weight as a judge only from audited claims, which its ratings of others cannot affect. Following the bipartite assignment of Lev et al.\ \cite{lev2023peernomination} and partition mechanisms \cite{alon2011sum}, ratings will flow only between groups. Agents that share sources or an underlying model go in the same group, since a model tends to recognize and favor its own outputs \cite{panickssery2024recognize}. Existing guarantees cover one agent deviating at a time, so we will measure how much a coalition of $k$ agents can still gain and flag groups whose agreement exceeds what their measured dependence explains. For selecting top candidates, we will adapt randomized impartial selection \cite{kurokawa2015impartial}, treating RG1's skill estimates as predictions that improve selection without breaking impartiality \cite{cembrano2025predictions}, and report the accuracy that impartiality costs. These mechanisms have been evaluated mostly on simulated reviewers; language-model reviewers allow tests at scale, including on scientific peer review.

\rg{RG2. Decide what to check, and settle disputes with evidence}
Both tasks follow one rule: \emph{a challenge can change the group's answer only through verifiable evidence}, meaning a record or test result obtained by a check, a quotation verified against its source, or a reasoning step simple enough to verify directly. Confidence, persistence, and the number of voices do not count.

\task{Task 2.1: Decide what to check.} Checking every claim is too costly, and checking only where agents disagree misses the most dangerous errors, because correlated agents agree on the same wrong answer and colluding agents can simply agree. We will treat the choice of which claims to check as a game against an adversary that knows our policy but not its random draws. Each batch of checks mixes four streams: a uniform random sample, which keeps estimates unbiased and leaves no claim safe from checking; claims where the weighted disagreement could change the decision, found by aligning the agents' claims and ranking suspect steps with FG-PRM \cite{li2025fgprm}; claims on which agents that RG1 estimates to be dependent agree more than their shared evidence explains; and claims flagged by the dishonesty signals of Task 2.2. Some trusted checks are provably necessary when honest raters may be a minority \cite{steinhardt2016avoiding}, and spot checks against occasional trusted evaluations give stronger incentives than many peer-prediction mechanisms while using less ground truth \cite{gao2016incentivizing}. We will derive how many checks are needed to expose a coalition of $k$ agents that behaves well whenever it expects a check, and how to split a budget across the four streams. In the vessel example, one check of the latest port-authority entry settles the dispute, and random checks of earlier claims can catch the compromised agent where it did not expect a check. Baselines check only disagreements, only the least confident claims, or claims chosen uniformly at random, at matched budgets.

\task{Task 2.2: Debate among many agents, settled by evidence.} Some disputes no single check can settle: the solutions differ too much to align, or no record, tool, or connection is available. For these we will extend debate to many agents of unknown and correlated reliability, comparing adversarial debate before a judge \cite{irving2018debate,khan2024debating} with collaborative disagreement resolution, in which agents locate the crux of their disagreement and revise \cite{jiang2026disagreement}. Each side is represented by its most reliable agents from RG1, and agents whose errors are strongly correlated share one speaker, so repeating a view adds no weight \cite{estornell2024debate}. Discussion narrows the dispute to a single claim; if that claim becomes checkable, Task 2.1 checks it. Only verified quotations and directly verifiable steps count as evidence, and the opposing side may add missing context, because true fragments can still mislead \cite{hu2026lying}. A panel of agents that took no side judges the dispute, weighted by RG1, with a designated judge model as a fallback; no agent judges its own answer. We will track three signs of dishonesty (arguing against one's own sealed first answer without new evidence, citing evidence that does not support the claim, and changing others' answers without evidence) and feed them into RG1 and Task 2.1. The evaluation measures how often a wrong answer is corrected, how often a correct one is lost, and how many dishonest agents each format tolerates.

\rg{RG3. Learn from verified outcomes}
\task{Task 3.1: Update trust without bias.} Each dispute settled by a check becomes a new audited claim that updates RG1. Targeted checks are chosen because something looked wrong, so they over-represent errors; we will estimate accuracy from the random stream with prediction-powered inference \cite{angelopoulos2023ppi} and weight targeted audits by their known selection probabilities. Outcomes decided only by debate update the weights provisionally and never become audits, so persuasive errors cannot reinforce themselves. Trust is tracked per question type and decays over time, which limits how much an agent gains by behaving well early and defecting later. We will also learn which kind of check resolves which kind of dispute, extending HiPRAG's rewards for searching only when needed \cite{wu2026hiprag}.

\task{Task 3.2: Reward verified contributions.} Training debaters to win can improve judge accuracy \cite{arnesen2024selfplay}, but rewarding persuasion has also taught models to exploit the judge, including by fabricating quotations \cite{kenton2026debate}, and to mislead human evaluators \cite{wen2025mislead}. We will reward each message for its verified effect. A message earns reward when it introduces verified evidence that moves the group toward the correct answer, or when its author revises in response to such evidence; changing others' answers without evidence, giving way without evidence, and abandoning a correct first answer are penalized. Deciding which message moved the group is a credit-assignment problem, which we approach with FG-PRM's step-level rewards and the PI's process-guided fine-tuning, building on multi-agent reinforcement learning for language models \cite{park2025maporl,subramaniam2025multiagent}. Shared rewards between agents that check each other can produce collusion \cite{shi2026collusion}, so each agent is rewarded separately and training episodes are audited at random. The main experiment compares training to win with training to contribute, measuring accuracy, lost correct answers, fabricated evidence, and emergent collusion.

\rg{Task 4 (shared). Testbeds and a red-team evaluation}
All scenarios are synthetic or drawn from public data. (A) Math and code problems solved by agents of known skill, some given the same flawed hint, so expertise and dependence are known by construction. (B) Multi-document report assessment built with our MEQA question-generation pipeline \cite{li2024meqa} and from MuSiQue \cite{trivedi2022musique}, including synthetic maritime scenarios with outdated and conflicting reports. (C) Scientific claims and peer review, using public conference reviews and author rebuttals, with claims later confirmed or contradicted by reproducibility studies serving as audits. On these testbeds we will run a control evaluation \cite{greenblatt2024control}: a red team trains agents to collude, to coordinate through stylistic signals, and to attack only when they expect no check, while the blue team uses the methods above, trusted weaker models, canary cases with known answers, and limited audits. We will also extend a public two-agent collusion environment \cite{shi2026collusion} to many agents. Results are curves of safety (wrong decisions avoided) against usefulness (correct decisions made without escalation) as the audit budget varies, compared with majority vote, Dawid-Skene weighting, PRD, and plain and two-sided debate at matched compute and checking budgets.

\section*{5\quad Limitations of our approach}
If every agent shares the same misconception and nothing can be checked, no aggregation or debate can recover the truth; our methods should report such claims as unresolved. The approach needs some trusted checks, and checks can themselves be wrong, so we will measure how sensitive each method is to unreliable checks and judges. Robustness will be shown only against the dishonest and colluding behaviors we test, because agents can coordinate through hidden signals that monitoring and paraphrasing do not fully remove \cite{mathew2025hidden}. Operational naval use would require domain data and expert validation beyond this fundamental research.

\section*{6\quad Impact if successful}
The project will show when agreement among AI agents counts as evidence and how many checks it takes to make it count. Each task targets a property of naval settings: unknown reliability (Tasks 1.1 and 3.1), late and sparse ground truth (Tasks 2.1 and 3.1), compromised agents (Tasks 1.2, 2.2, and 3.2), and costly or unavailable checks (Tasks 2.1 and 2.2). For an analyst using several AI tools, the result would show which tools' judgments carry weight and why, identify the claim still in dispute, and say when a decision needs human review. We will release code, testbeds, and evaluation protocols.

\clearpage
\begingroup
\small
\setlength{\parskip}{0pt}
\begin{thebibliography}{99}
\setlength{\itemsep}{2.5pt}
\bibitem{zheng2023judging} L. Zheng et al. (2023). \href{https://proceedings.neurips.cc/paper_files/paper/2023/hash/91f18a1287b398d378ef22505bf41832-Abstract.html}{Judging LLM-as-a-Judge with MT-Bench and Chatbot Arena}. NeurIPS 36, Datasets and Benchmarks, pp. 46595--46623.
\bibitem{verga2024juries} P. Verga et al. (2024). \href{https://arxiv.org/abs/2404.18796}{Replacing Judges with Juries: Evaluating LLM Generations with a Panel of Diverse Models}. arXiv:2404.18796.
\bibitem{kim2025correlated} E. M. Kim, A. Garg, K. Peng, and N. Garg (2025). \href{https://proceedings.mlr.press/v267/kim25e.html}{Correlated Errors in Large Language Models}. ICML, PMLR 267:30038--30066.
\bibitem{goel2025great} S. Goel, J. Str\"uber, I. A. Auzina, K. K. Chandra, P. Kumaraguru, D. Kiela, A. Prabhu, M. Bethge, and J. Geiping (2025). \href{https://proceedings.mlr.press/v267/goel25b.html}{Great Models Think Alike and this Undermines AI Oversight}. ICML, PMLR 267:19621--19678.
\bibitem{hossain2026agreement} E. Hossain, N. Yousefi, and S.-N. Lim (2026). \href{https://arxiv.org/abs/2609.22512}{Agreement Overstates Evidence: Error Dependence in LLM Judge Consensus}. arXiv:2609.22512.
\bibitem{kohli2026nine} G. Kohli (2026). \href{https://arxiv.org/abs/2605.29800}{Nine Judges, Two Effective Votes: Correlated Errors Undermine LLM Evaluation Panels}. arXiv:2605.29800.
\bibitem{shu2026unanimity} Y. Shu (2026). \href{https://arxiv.org/abs/2609.26145}{Unanimity Without Persuasion: A Single Round of Debate Erases the Disagreement That Verification Needs}. arXiv:2609.26145.
\bibitem{shi2026collusion} X. Shi, Y. Zhang, and D. Yang (2026). \href{https://arxiv.org/abs/2609.24967}{Emergent Collusion in Long-Horizon LLM Agent Interaction}. arXiv:2609.24967.
\bibitem{dawid1979maximum} A. P. Dawid and A. M. Skene (1979). \href{https://doi.org/10.2307/2346806}{Maximum Likelihood Estimation of Observer Error-Rates Using the EM Algorithm}. Journal of the Royal Statistical Society, Series C (Applied Statistics) 28(1):20--28.
\bibitem{whitehill2009whose} J. Whitehill, P. Ruvolo, T. Wu, J. Bergsma, and J. Movellan (2009). \href{https://proceedings.neurips.cc/paper/2009/hash/f899139df5e1059396431415e770c6dd-Abstract.html}{Whose Vote Should Count More: Optimal Integration of Labels from Labelers of Unknown Expertise}. NeurIPS 22, pp. 2035--2043.
\bibitem{dong2009dependence} X. L. Dong, L. Berti-Equille, and D. Srivastava (2009). \href{https://doi.org/10.14778/1687627.1687689}{Integrating Conflicting Data: The Role of Source Dependence}. PVLDB 2(1):550--561.
\bibitem{jaffe2016dependent} A. Jaffe, E. Fetaya, B. Nadler, T. Jiang, and Y. Kluger (2016). \href{https://proceedings.mlr.press/v51/jaffe16.html}{Unsupervised Ensemble Learning with Dependent Classifiers}. AISTATS, PMLR 51:351--360.
\bibitem{zhu2026stopping} B. Zhu, Y. Xie, and Y. Rao (2026). \href{https://arxiv.org/abs/2608.19802}{Stopping and Routing LLM Judge Panels}. arXiv:2608.19802.
\bibitem{kurokawa2015impartial} D. Kurokawa, O. Lev, J. Morgenstern, and A. D. Procaccia (2015). \href{https://www.ijcai.org/Abstract/15/088}{Impartial Peer Review}. IJCAI, pp. 582--588.
\bibitem{alon2011sum} N. Alon, F. Fischer, A. D. Procaccia, and M. Tennenholtz (2011). \href{https://doi.org/10.1145/2000378.2000390}{Sum of Us: Strategyproof Selection from the Selectors}. TARK, pp. 101--110.
\bibitem{flanigan2021fair} B. Flanigan, P. G\"olz, A. Gupta, B. Hennig, and A. D. Procaccia (2021). \href{https://doi.org/10.1038/s41586-021-03788-6}{Fair Algorithms for Selecting Citizens' Assemblies}. Nature 596:548--552.
\bibitem{flanigan2024manipulation} B. Flanigan, J. Liang, A. D. Procaccia, and S. Wang (2024). \href{https://doi.org/10.1609/aaai.v38i9.28827}{Manipulation-Robust Selection of Citizens' Assemblies}. AAAI 38(9):9696--9703.
\bibitem{khalaf2026maximal} H. Khalaf, S. L. Wang, D. Halpern, I. Shapira, F. P. Calmon, and A. D. Procaccia (2026). \href{https://arxiv.org/abs/2602.21297}{Robust AI Evaluation through Maximal Lotteries}. ICML.
\bibitem{lev2023peernomination} O. Lev, N. Mattei, P. Turrini, and S. Zhydkov (2023). \href{https://doi.org/10.1016/j.artint.2022.103843}{PeerNomination: A Novel Peer Selection Algorithm to Handle Strategic and Noisy Assessments}. Artificial Intelligence 316:103843.
\bibitem{du2024debate} Y. Du, S. Li, A. Torralba, J. B. Tenenbaum, and I. Mordatch (2024). \href{https://proceedings.mlr.press/v235/du24e.html}{Improving Factuality and Reasoning in Language Models through Multiagent Debate}. ICML, PMLR 235:11733--11763.
\bibitem{choi2025debate} H. K. Choi, X. Zhu, and S. Li (2025). \href{https://arxiv.org/abs/2508.17536}{Debate or Vote: Which Yields Better Decisions in Multi-Agent Large Language Models?} NeurIPS.
\bibitem{sharma2024sycophancy} M. Sharma et al. (2024). \href{https://proceedings.iclr.cc/paper_files/paper/2024/hash/0105f7972202c1d4fb817da9f21a9663-Abstract-Conference.html}{Towards Understanding Sycophancy in Language Models}. ICLR.
\bibitem{laban2023flipflop} P. Laban, L. Murakhovs'ka, C. Xiong, and C.-S. Wu (2023). \href{https://arxiv.org/abs/2311.08596}{Are You Sure? Challenging LLMs Leads to Performance Drops in the FlipFlop Experiment}. arXiv:2311.08596.
\bibitem{amayuelas2024attack} A. Amayuelas, X. Yang, A. Antoniades, W. Hua, L. Pan, and W. Y. Wang (2024). \href{https://arxiv.org/abs/2406.14711}{MultiAgent Collaboration Attack: Investigating Adversarial Attacks in Large Language Model Collaborations via Debate}. Findings of EMNLP, pp. 6929--6948.
\bibitem{irving2018debate} G. Irving, P. Christiano, and D. Amodei (2018). \href{https://arxiv.org/abs/1805.00899}{AI Safety via Debate}. arXiv:1805.00899.
\bibitem{khan2024debating} A. Khan et al. (2024). \href{https://proceedings.mlr.press/v235/khan24a.html}{Debating with More Persuasive LLMs Leads to More Truthful Answers}. ICML, PMLR 235:23662--23733.
\bibitem{jiang2026disagreement} Y. Jiang, C. Chen, T. Wu, L. Sun, H. Liu, S. Feng, and C. Tan (2026). \href{https://arxiv.org/abs/2607.01251}{Collaborative Disagreement Resolution for Scalable Oversight}. ICML.
\bibitem{hu2026lying} J. Hu, X. Huang, Y. Sun, Y. Dong, and X. Huang (2026). \href{https://arxiv.org/abs/2601.01685}{Lying with Truths: Open-Channel Multi-Agent Collusion for Belief Manipulation via Generative Montage}. ACL.
\bibitem{qiu2026truthfulness} T. A. Qiu, M. Carroll, and C. Allen (2026). \href{https://arxiv.org/abs/2601.20299}{Truthfulness Despite Weak Supervision: Evaluating and Training LLMs Using Peer Prediction}. ICLR.
\bibitem{greenblatt2024control} R. Greenblatt, B. Shlegeris, K. Sachan, and F. Roger (2024). \href{https://arxiv.org/abs/2312.06942}{AI Control: Improving Safety Despite Intentional Subversion}. ICML, PMLR 235:16295--16336.
\bibitem{li2024prd} R. Li, T. Patel, and X. Du (2024). \href{https://openreview.net/forum?id=YVD1QqWRaj}{PRD: Peer Rank and Discussion Improve Large Language Model based Evaluations}. TMLR.
\bibitem{li2025fgprm} R. Li, Z. Luo, and X. Du (2025). \href{https://aclanthology.org/2025.findings-emnlp.228/}{FG-PRM: Fine-grained Hallucination Detection and Mitigation in Language Model Mathematical Reasoning}. Findings of EMNLP, pp. 4247--4278.
\bibitem{wu2026hiprag} P. Wu, M. Zhang, K. Wan, W. Zhao, K. He, X. Du, and Z. Chen (2026). \href{https://proceedings.iclr.cc/paper_files/paper/2026/hash/85255cf8f2d2318cc7cffd0b1bc66625-Abstract-Conference.html}{HiPRAG: Hierarchical Process Rewards for Efficient Agentic Retrieval Augmented Generation}. ICLR.
\bibitem{panickssery2024recognize} A. Panickssery, S. R. Bowman, and S. Feng (2024). \href{https://arxiv.org/abs/2404.13076}{LLM Evaluators Recognize and Favor Their Own Generations}. NeurIPS 37.
\bibitem{cembrano2025predictions} J. Cembrano, F. Fischer, and M. Klimm (2025). \href{https://arxiv.org/abs/2510.19002}{Impartial Selection with Predictions}. NeurIPS.
\bibitem{steinhardt2016avoiding} J. Steinhardt, G. Valiant, and M. Charikar (2016). \href{https://arxiv.org/abs/1606.05374}{Avoiding Imposters and Delinquents: Adversarial Crowdsourcing and Peer Prediction}. NeurIPS 29.
\bibitem{gao2016incentivizing} A. Gao, J. R. Wright, and K. Leyton-Brown (2016). \href{https://arxiv.org/abs/1606.07042}{Incentivizing Evaluation via Limited Access to Ground Truth: Peer-Prediction Makes Things Worse}. arXiv:1606.07042.
\bibitem{estornell2024debate} A. Estornell and Y. Liu (2024). \href{https://proceedings.nips.cc/paper_files/paper/2024/file/32e07a110c6c6acf1afbf2bf82b614ad-Paper-Conference.pdf}{Multi-LLM Debate: Framework, Principals, and Interventions}. NeurIPS 37.
\bibitem{angelopoulos2023ppi} A. N. Angelopoulos, S. Bates, C. Fannjiang, M. I. Jordan, and T. Zrnic (2023). \href{https://arxiv.org/abs/2301.09633}{Prediction-Powered Inference}. Science.
\bibitem{arnesen2024selfplay} S. Arnesen, D. Rein, and J. Michael (2024). \href{https://arxiv.org/abs/2409.16636}{Training Language Models to Win Debates with Self-Play Improves Judge Accuracy}. arXiv:2409.16636.
\bibitem{kenton2026debate} Z. Kenton, L. Janzer, R. Greig, T. H. Teh, K. Tyshchuk, J. Brown-Cohen, H. Edwards, S. Rajamanoharan, N. Y. Siegel, N. Jaques, and R. Shah (2026). \href{https://arxiv.org/abs/2608.17776}{Debate Training Reduces Reward Hacking in RLAIF}. arXiv:2608.17776.
\bibitem{wen2025mislead} J. Wen, R. Zhong, A. Khan, E. Perez, J. Steinhardt, M. Huang, S. R. Bowman, H. He, and S. Feng (2025). \href{https://arxiv.org/abs/2409.12822}{Language Models Learn to Mislead Humans via RLHF}. ICLR.
\bibitem{park2025maporl} C. Park, S. Han, X. Guo, A. E. Ozdaglar, K. Zhang, and J.-K. Kim (2025). \href{https://aclanthology.org/2025.acl-long.1459}{MAPoRL: Multi-Agent Post-Co-Training for Collaborative Large Language Models with Reinforcement Learning}. ACL, pp. 30215--30248.
\bibitem{subramaniam2025multiagent} V. Subramaniam, Y. Du, J. B. Tenenbaum, A. Torralba, S. Li, and I. Mordatch (2025). \href{https://proceedings.iclr.cc/paper_files/paper/2025/hash/1d50221b03ced7152d8750c9d1e66cba-Abstract-Conference.html}{Multiagent Finetuning: Self Improvement with Diverse Reasoning Chains}. ICLR.
\bibitem{li2024meqa} R. Li, Z. Wang, S. Q. Tran, L. Xia, and X. Du (2024). \href{https://proceedings.neurips.cc/paper_files/paper/2024/hash/e560a0b22e4432003d0dba63ff8dc457-Abstract.html}{MEQA: A Benchmark for Multi-hop Event-centric Question Answering with Explanations}. NeurIPS 37, Datasets and Benchmarks, pp. 126835--126862.
\bibitem{trivedi2022musique} H. Trivedi, N. Balasubramanian, T. Khot, and A. Sabharwal (2022). MuSiQue: Multihop Questions via Single-hop Question Composition. TACL 10:539--554.
\bibitem{mathew2025hidden} Y. Mathew, O. Matthews, R. McCarthy, J. Velja, C. Schroeder de Witt, D. Cope, and N. Schoots (2025). \href{https://aclanthology.org/2025.ijcnlp-long.34/}{Hidden in Plain Text: Emergence and Mitigation of Steganographic Collusion in LLMs}. IJCNLP-AACL, pp. 585--624.
\end{thebibliography}
\endgroup
\end{document}
